\documentclass[conference]{IEEEtran}
\usepackage[T1]{fontenc}
\usepackage{amsmath,amssymb,graphicx,booktabs,url,balance}
\newtheorem{proposition}{Proposition}
\newcommand{\R}{\mathbb{R}}
\newcommand{\norm}[1]{\left\lVert#1\right\rVert}
\title{Boundary-Aware Static Robot Deployment:\\Sampled Safety and Limits of Transport Preparation}
\author{\IEEEauthorblockN{Development manuscript}\IEEEauthorblockA{Author information pending}}
\begin{document}
\maketitle
\begin{abstract}
Placing robots around an object is a distinct problem from reliably transporting
it. We study boundary-aware deployment through limited-range centroidal targets
and a velocity safety filter, separating an exact static oracle-map analysis
from locally sensed implementations and downstream contact experiments. For a
fixed polygon and single-integrator robots, a feature-cover construction gives
whole-hold safety under explicit initialization, communication, sampling and
projection assumptions. A projected-residual bound clarifies why low commanded
motion need not imply complete coverage. In a frozen 169-episode campaign, a
local segment candidate passes nine of nine static path checks, versus one of
nine for the original point-sample implementation, while reducing the measured
deployment cost. Its gap guard intervenes on 5,315 of 5,400 frames. Transport
completion remains limited: boundary CVT completes 5/12 new-seed tasks and 1/18
perturbed tasks. An independent contact-engine case also fails. These results
support a restricted static deployment claim and motivate, rather than establish,
reliable transport with unknown moving boundaries.
\end{abstract}
\begin{IEEEkeywords}
Multi-robot deployment, coverage control, sampled safety, cooperative transport.
\end{IEEEkeywords}

\section{Introduction}
Robots that must push an unfamiliar object first need useful positions around
its boundary. A team can approach the boundary quickly while leaving important
regions unsupported; a configuration can also appear stationary because safety
constraints block its nominal motion. Once contact starts, estimation drift,
asymmetric forces and release timing introduce further failure modes. A single
success label obscures these distinctions.

We investigate a boundary-aware deployment controller with separate information
permissions for control and evaluation. Limited-range Voronoi cells allocate
observed boundary mass, while a safety filter modifies the resulting velocity.
The theoretical reference uses a fixed known polygon. The local implementation
receives only sensed or relayed boundary information; true geometry is reserved
for an external evaluator. Contact transport is evaluated as a downstream task,
without extending the static theorem to moving objects.

The contributions are an explicit static feature-cover safety argument, a
residual interpretation for the saturated-then-projected centroidal law, and an
auditable comparison that retains failed and censored episodes. The empirical
segment repair improves the observed static paths with essentially unchanged
cost reduction, but its conservatism and incomplete geometry prevent a local-map
safety theorem. Preparation and transport results are reported separately.
We do not claim general transport superiority or engine invariance.

\section{Related Work}
Voronoi-based coverage control connects locational costs to distributed
centroidal motion~\cite{cortes2004}. We use that connection to allocate boundary
density, with truncated cells and a fixed-density theoretical reference. Safety
projection and numerical centroid errors change the interpretation of the
classical centroidal equilibrium; coverage convergence cannot simply be inherited.

Control barrier functions express forward-invariance requirements as constraints
on admissible controls and can be combined with performance objectives in a
quadratic program~\cite{ames2017}. Our concern is the held command between samples,
including switches of the closest polygon feature. A feasible instantaneous
inequality against an incomplete observed boundary is insufficient.

Caging-based decentralized manipulation maintains object closure while moving
toward a goal~\cite{pereira2004}. Sequential caging has also addressed unknown
object shapes, distributing contact locations before transport and evaluating
both simulation and physical robots~\cite{vardharajan2021}. Those results make
unknown-shape transport an established research direction, not a novelty claim
of the present work. Here the narrower question is what a boundary-density
deployment objective and its static safety filter actually establish before
dynamic contact. We do not equate proximity to a boundary with object closure
or sufficient contact wrench authority.

\section{Model, Information and Objectives}
Let $D\subset\R^2$ be a closed rectangle and $S\subset D$ a fixed simple polygon.
There are $N=16$ robots with $\dot p_i=u_i$, $\norm{u_i}\le u_{\max}$.
At $t_k=k\Delta$, all controls are computed from the same frozen configuration
$P_k$, then held for $\tau\in[0,\Delta]$. The theoretical oracle exposes the
polygon features. Local arms instead construct their maps from range observations
and one-hop communication. An independent evaluator may access $S$ and simulator
object states, but those values do not select local control commands.

For the analysis only, fix a nonnegative reference density $\phi^\star$ around a
uniform boundary offset. Direction-dependent, exploration and changing-map terms
are excluded. With $\Omega_i=V_i(P)\cap B(p_i,R)$, define
\begin{equation}
 H^\star(P)=\int_D \min\{R^2,\min_i\norm{q-p_i}^2\}\phi^\star(q)\,dq.
\end{equation}
For each positive-mass cell, $m_i^\star=\int_{\Omega_i}\phi^\star$ and
$c_i^\star=(m_i^\star)^{-1}\int_{\Omega_i}q\phi^\star(q)\,dq$.
The implementation computes $\hat m_i,\hat c_i$, projects an out-of-disk centroid
onto the local disk and uses $p_i$ when its numerical mass is below the floor.
The command is
\begin{equation}\label{eq:cascade}
 v_i=\operatorname{sat}_{u_{\max}}(k_c(\hat c_i-p_i)),\qquad
 U_i=\Pi_{F_{i,k}}(v_i).
\end{equation}
Here $F_{i,k}$ is the intersection of the speed ball, shared-responsibility
robot rows, rectangular wall rows and object-feature rows. The exact analysis
uses the original robust margin, with no infeasibility scaling or clamping.

Static quality is measured by cost reduction and displacement as well as safety.
In transport, controller CONTACT\_READY, a dwell of positive forward physical
contacts, G500, settling and final task completion are distinct events. Completion
requires the existing settling rule and final Euclidean position error at most
the task tolerance, $0.12$ m in the reported transport cases. An along-track
arrival estimate alone does not satisfy this definition.

\section{Static Sampled Safety and Residual Interpretation}
For a polygon feature $F$ (a closed edge segment), set
$h_F(p)=\operatorname{dist}(p,F)-r_{\mathrm{safe}}$. Include every feature in
\begin{equation}
\begin{split}
 \Phi(p)=\{F:\ &\operatorname{dist}(p,F)\le d_\partial(p)+2u_{\max}\Delta,\\
 &\operatorname{dist}(p,F)\le R_{\mathrm{row}}\},
\end{split}
\end{equation}
where $d_\partial$ is distance to the polygon boundary. Impose
\begin{equation}\label{eq:row}
 \nabla h_F(p)^T u\ge-\gamma_{\mathrm{obj}}h_F(p)+\rho.
\end{equation}
Distances are to finite segments, not their infinite supporting lines.

\subsection{Assumptions and whole-hold argument}
The initial configuration lies in $D$, has pair separation at least $d_{\min}$,
and lies outside $S$ with signed clearance at least
$r_{\mathrm{safe}}+\rho/\gamma_{\mathrm{obj}}$. Geometry and density are fixed;
communication has no dropout and range at least $2R$, including all pairs within
$d_{\min}+2u_{\max}\Delta$. The robot rows use the same frozen positions and
symmetric responsibility. Require $\gamma_{\mathrm{obj}}\Delta\le1$,
$\gamma_{\mathrm{agent}}\Delta\le1$, and
\begin{equation}\label{eq:range}
 R_{\mathrm{row}}-u_{\max}\Delta\ge
 r_{\mathrm{safe}}+\rho/\gamma_{\mathrm{obj}}.
\end{equation}
All projections and geometry in the following argument use exact arithmetic.

\begin{proposition}[Static safety under the stated assumptions]
The exact cascade preserves object clearance
$r_{\mathrm{safe}}+\rho/\gamma_{\mathrm{obj}}$, pair separation $d_{\min}$ and
membership in $D$ throughout every hold. Its feasible velocity sets contain zero
at every sample, so their Euclidean projections exist recursively.
\end{proposition}
\begin{IEEEproof}
Distance to a closed convex segment is convex. For each included feature,
\eqref{eq:row} and the supporting-gradient inequality give
\begin{equation}\label{eq:hold}
 h_F(p+\tau u)\ge(1-\gamma_{\mathrm{obj}}\tau)h_F(p)+\rho\tau.
\end{equation}
An in-range feature excluded by the cover is initially more than
$2u_{\max}\Delta$ beyond the nearest boundary distance. Its distance decreases
by at most $u_{\max}\Delta$, by the Lipschitz property. Range-excluded features
remain beyond the required margin by \eqref{eq:range}. Thus the minimum boundary
distance retains the margin even at feature switches or when the included set
is empty. Starting outside, continuity excludes a first boundary crossing.
For included pairs, the sum of the two responsibility inequalities bounds the
initial relative-distance derivative; convexity of squared distance and
$\gamma_{\mathrm{agent}}\Delta\le1$ preserve separation. Omitted pairs cannot
close the initial excess distance within a hold. Rectangular wall constraints
preserve each coordinate throughout the segment. These margins make zero feasible
in every stacked row at the next sample. Induction proves the assertion.
\end{IEEEproof}

The reference parameters are $\Delta=0.05$ s, $k_c=0.9$, $u_{\max}=0.35$ m/s,
$R=0.8$ m, $R_{\mathrm{row}}=0.60$ m, $d_{\min}=0.28$ m,
$r_{\mathrm{safe}}=0.065$ m, $\rho=0.02$, $\gamma_{\mathrm{obj}}=8$ s$^{-1}$
and $\gamma_{\mathrm{agent}}=6$ s$^{-1}$. The object margin is $67.5$ mm;
the local empirical path acceptance threshold is $65$ mm. These are different
quantities. Full derivations and parameter ledgers accompany the source.

\subsection{Cost dissipation does not imply complete coverage}
Define $J_k=\sum_i m_{i,k}^\star\norm{U_{i,k}}^2$ and
$E_k=\sum_i m_{i,k}^\star\norm{\hat c_{i,k}-c_{i,k}^\star}^2$, with zero-mass
cells contributing zero. Freezing the current partition gives a quadratic
majorant of $H^\star$; reassigning points to their nearest truncated cell only
reduces that majorant. Projection optimality with zero feasible then yields,
for $a=2/k_c-\Delta>0$,
\begin{equation}\label{eq:dissipation}
 \overline J\le \frac{2H^\star(P_0)}{aK\Delta}+\frac{4\overline E}{a^2}.
\end{equation}
To see this, write $v_i=\lambda_i(\hat c_i-p_i)$, with
$0<\lambda_i\le k_c$. Projection optimality and zero feasibility imply
$v_i^TU_i\ge\norm{U_i}^2$. The partition majorant then gives
$H^\star(P_{k+1})-H^\star(P_k)\le-\Delta aJ_k+2\Delta\sqrt{E_kJ_k}$.
Young's inequality bounds the right-hand side by
$-\Delta aJ_k/2+2\Delta E_k/a$; summing and using $H^\star\ge0$
proves \eqref{eq:dissipation}. This argument does not assume numerical
monotonicity of the observed cost.

To clarify its meaning, let $r_i$ be \eqref{eq:cascade} computed with $c_i^\star$
in place of $\hat c_i$ and the same $F_{i,k}$. Nonexpansiveness of saturation
and projection gives $\norm{r_i-U_i}\le k_c\norm{c_i^\star-\hat c_i}$. Hence
\begin{equation}
 G_k:=\sum_i m_{i,k}^\star\norm{r_i}^2\le2J_k+2k_c^2E_k,
\end{equation}
and \eqref{eq:dissipation} bounds the average exact-centroid projected residual.
Zero residual means no descent through the filter's frozen feasible velocity
set. It does not imply stationarity of the original geometric-position problem
without an additional tangent-cone equivalence. Nearly massless cells can also
hide poor spatial support. Therefore neither a small $J$ nor a small $G$
establishes full boundary coverage, useful contact capacity or transport success.

\subsection{Local maps, arithmetic and moving boundaries}
A sufficient local-map extension would require certified coverage of every true
boundary point relevant to the hold tube within error $\epsilon$, with omitted
features independently excluded. A whole-hold estimated clearance of at least
$r_{\mathrm{safe}}+\epsilon$ would then transfer to true clearance by the triangle
inequality. Our sensor does not provide this coverage certificate. In particular,
clipping a measured observation gap does not certify missing geometry.

Floating-point row residuals also require a separate budget. If the physical
inequality has violation at most $\epsilon_u$, the last term in \eqref{eq:hold}
becomes $(\rho-\epsilon_u)\tau$. Row coefficient errors, state error and tracking
error contribute in addition to solver tolerance. Thus a floating-point
\texttt{optimal} status is not an exact-arithmetic certificate. Moving polygons
would further require swept-feature bounds, uncertain boundary velocities and
recursive feasibility under contact dynamics. Those extensions remain open.

\section{Experiments and Failure Evidence}
\subsection{Protocol and provenance}
The frozen draft-v2 campaign contains 169 planned episodes: 27 static comparisons,
27 development transport comparisons, 15 target ablations, six active-exploration
diagnostics, 36 generalization tasks, 54 perturbation tasks, two QP references
and two contact-engine references. All planned cases remain in their denominators.
The source snapshot, expanded configurations, initial-state identities and product
hashes are archived. The implementation audit passes with zero execution errors;
that gate does not require scientific task success. Two serial repeats, excluded
from scientific denominators, reproduce their parallel results.

Static comparisons use rectangle, L and C polygons and seeds 2, 5 and 8, with
600 complete control holds each. The oracle, original local samples and local
segment candidate share the task and initial state. Transport compares boundary
CVT, nearest observed targets and local spacing; only the pre-CONTACT\_READY
deployment target differs. Generalization uses 12 new seeds with the frozen
layout/scale factors. Perturbations use two independent seeds repeated across
shapes and noise, communication loss and slower sensing. These are two seed
clusters, not 18 independent replications per method.

\begin{table}[t]
\caption{Frozen v2 static paths: nine cases per arm; required clearance 65 mm.}
\label{tab:static}\centering
\begin{tabular}{lrrr}\toprule
Arm & Safe paths & Min. (mm) & Mean $H_f/H_0$\\\midrule
Original local samples & 1/9 & 48.33 & 0.1533\\
Segment candidate & 9/9 & 66.04 & 0.1544\\
Oracle reference & 9/9 & 67.50 & 0.1534\\\bottomrule
\end{tabular}
\end{table}

\subsection{Static deployment}
Table~\ref{tab:static} shows that the local segment candidate passes all nine
whole-hold path checks; the original local point formulation violates eight.
The segment candidate's smallest observed margin over the threshold is only
$1.04$ mm. Mean robot displacement per episode ranges from $0.810$ to $1.013$ m,
so the cost reduction does not result from remaining stationary. However, the
gap guard truncates 5,315 of 5,400 frames. This frequent intervention and the
limited three-shape geometry coverage prevent broad generalization.

The independent planar-QP and CVXPY references both pass their 600-frame checks,
with maximum recorded position difference $1.82$ mm. We report agreement of the
checked safety outcomes, not equality of numerical trajectories. Finite path
verification remains distinct from the exact proposition's assumptions.

\subsection{Preparation versus final transport}
\begin{table}[t]
\caption{Frozen v2 task completion, including failed and censored tasks.}
\label{tab:transport}\centering
\begin{tabular}{lrrr}\toprule
Population & CVT & Nearest & Spacing\\\midrule
Development & 5/9 & 5/9 & 4/9\\
New-seed generalization & 5/12 & 3/12 & 1/12\\
Perturbation robustness & 1/18 & 2/18 & 1/18\\\bottomrule
\end{tabular}
\end{table}
In generalization, CVT reaches controller readiness $1.364$ s earlier on average
than nearest-target deployment over 11 complete pairs; one pair is censored. The
exploratory seed-cluster bootstrap interval for CVT minus nearest is
$[-2.736,-0.145]$ s. Against local spacing, the difference is $-1.632$ s with
interval $[-6.668,1.582]$ s. Complete-pair selection can bias readiness differences.
Neither comparison establishes transport reliability; the actual task rates
are given in Table~\ref{tab:transport}. Development and robustness have only
three and two independent seed clusters, respectively, so their rates are
descriptive and receive no inferential interval.

The main exploration gain is zero, and removing it produces identical replay
arrays in all three baseline shapes. It is a null intervention, not evidence
for the contribution of exploration. A separate gain-six diagnostic has mixed
outcomes. One seed per shape is insufficient for an ablation benefit claim.
In the independent-engine pair, penalty contact completes at $0.111$ m endpoint
error; Pymunk reaches $0.378$ m and triggers its watchdog. Successful recording
and execution in both engines do not establish physical-model invariance.

\subsection{Failure classification and reliability repairs}
Reanalysis of the 140 transport receipts retains the frozen outcomes and uses
each task's original tolerance. For CVT development failures, three are dominated
by along-track undershoot and one by overshoot. For its 17 robustness failures,
eight exceed tolerance in the final transverse component, eight exceed it in
both transverse and longitudinal components, and one is longitudinal undershoot.
These are geometric outcome classes. Missing physical-contact dwell and missing
phase events are recorded separately and are not assigned as causal mechanisms
from compact receipts alone.

Development inspection also found that slow planning could retain a nominal
command across a supervisor change, and a robot whose progress estimate lagged
the team could still enter the push set after HOLD. The default revision
invalidates held commands when supervisor permissions change and explicitly
prohibits pushing in HOLD. More aggressive changes to brake-integral release
and feedback routing were tested but caused development regressions; they remain
experimental options disabled by default. This negative result cautions against
equating a plausible control-law change with improved task reliability.
No revision changes tolerances, horizons or the frozen results above.
Regression tests establish specific control responses; development replays are
separate from independent validation and do not establish robust transport.

\balance
\section{Limitations and Conclusion}
The strongest conclusion is static sampled deployment under explicit fixed-map
assumptions, supported by independent path observations for a local segment
candidate on nine cases. Static deployment quality and downstream controller
readiness provide useful measurements, but neither establishes caging or contact
wrench sufficiency. Robustness failures and the Pymunk discrepancy are material
limits, not execution errors to exclude.

Further validation must freeze additional geometry, occlusion and sampling
conditions before observing outcomes, report gap-guard intervention alongside
deployment quality, and use independent seed clusters for transport comparisons.
Mechanism studies need frame traces for estimate bias, contact-loss response,
brake entry and terminal error. Dynamic safety and hardware deployment require
additional modeling and sensing evidence. The present results delimit a static
deployment contribution while identifying the measurements needed to assess
reliable transport in subsequent work.

\begin{thebibliography}{4}
\bibitem{cortes2004} J. Cort\'es, S. Mart\'inez, T. Karatas, and F. Bullo,
``Coverage control for mobile sensing networks,'' \emph{IEEE Trans. Robot. Autom.},
vol. 20, no. 2, pp. 243--255, 2004, doi: 10.1109/TRA.2004.824698.
\bibitem{ames2017} A. D. Ames, X. Xu, J. W. Grizzle, and P. Tabuada,
``Control barrier function based quadratic programs for safety critical systems,''
\emph{IEEE Trans. Autom. Control}, vol. 62, no. 8, pp. 3861--3876, 2017,
doi: 10.1109/TAC.2016.2638961.
\bibitem{pereira2004} G. A. S. Pereira, M. F. M. Campos, and V. Kumar,
``Decentralized algorithms for multi-robot manipulation via caging,''
\emph{Int. J. Robot. Res.}, vol. 23, no. 7--8, 2004,
doi: 10.1177/0278364904045477.
\bibitem{vardharajan2021} V. S. Vardharajan, K. Soma, and G. Beltrame,
``Collective transport via sequential caging,'' arXiv:2106.03132, 2021.
\end{thebibliography}
\end{document}
